\documentclass[11pt]{article}
\usepackage[a4paper,margin=2.2cm]{geometry}
\usepackage{xcolor}
\usepackage{fontspec}
\usepackage{xeCJK}
\usepackage{enumitem}
\setmainfont{Times New Roman}
\setCJKmainfont{SimSun}
\setlist[itemize]{leftmargin=1.4em,itemsep=0.35em,topsep=0.35em}
\setlength{\parindent}{0pt}
\setlength{\parskip}{0.55em}
\pagestyle{plain}
\begin{document}
{\color[HTML]{1F5FD2}\LARGE\bfseries 融合教育政策要点解读\par}
{\color[HTML]{667085}\small 星轨原创教育资源：用于家庭与学校共同制定支持计划。\par}
\vspace{0.8em}
{\color[HTML]{111827}\large\bfseries 核心理念\par}
\begin{itemize}
\item 融合教育不是把儿童简单放进普通班级，而是用环境调整、课程弹性和同伴支持，让每个儿童都能有尊严地参与学习。
\item 判断一项支持是否有效，应看儿童是否更愿意参与、更能表达需要，并在真实课堂中逐步减少成人提示。
\end{itemize}
{\color[HTML]{111827}\large\bfseries 政策落地清单\par}
\begin{itemize}
\item 为儿童建立支持画像，记录兴趣优势、敏感情境、沟通方式、有效提示和合理便利措施。
\item 采用通用学习设计，在同一目标下提供图片、实物、口语说明、动作示范和同伴示范等多种进入方式。
\item 把支持嵌入日常流程，包括座位安排、活动转换、任务长度、反馈频率、休息方式和家校沟通。
\end{itemize}
{\color[HTML]{111827}\large\bfseries 团队协作\par}
\begin{itemize}
\item 每月至少一次由家长、班主任、特教教师和相关专业人员共同查看观察记录与作品样本。
\item 会议重点不是追责，而是确认哪些支持需要保留、淡化或加强，并写入下一阶段行动清单。
\end{itemize}
\vfill
{\color[HTML]{667085}\footnotesize 本材料为应用内原创教育内容，不能替代医疗、法律或正式评估建议。\par}
\end{document}
